% !TEX program = xelatex
% Lernplatt - by Can Siebert
% Kurs 71 (Dig 42) - Tag 2 - Loesungsheft
% UML, VSM & Process Mining als Dreiklang aus Struktur, Wertfluss und Datenrealitaet

\documentclass[11pt,a4paper,oneside]{article}

\usepackage{polyglossia}
\setdefaultlanguage{german}

\usepackage[a4paper,margin=2.5cm]{geometry}

\usepackage{fontspec}
\defaultfontfeatures{Ligatures=TeX}

\IfFontExistsTF{Charter}{%
  \setmainfont{Charter}%
}{%
  \IfFontExistsTF{Bitstream Charter}{%
    \setmainfont{Bitstream Charter}%
  }{%
    \setmainfont{TeX Gyre Schola}%
  }%
}

\IfFontExistsTF{Source Sans 3}{%
  \setsansfont{Source Sans 3}%
}{%
  \IfFontExistsTF{Source Sans Pro}{%
    \setsansfont{Source Sans Pro}%
  }{%
    \setsansfont{TeX Gyre Heros}%
  }%
}

\newcommand{\caveatfontfallback}{\itshape}
\IfFontExistsTF{Caveat}{%
  \newfontfamily\caveatfont{Caveat}[Scale=1.15]%
}{%
  \newcommand\caveatfont{\caveatfontfallback}%
}

\setmonofont{Latin Modern Mono}

\usepackage{microtype}
\linespread{1.15}

\usepackage{xcolor}
\definecolor{lpInk}{RGB}{20,28,38}
\definecolor{lpAccent}{RGB}{157,74,26}
\definecolor{lpAccentLight}{RGB}{246,228,212}
\definecolor{lpAmber}{RGB}{222,138,30}
\definecolor{lpAmberLight}{RGB}{255,243,217}
\definecolor{lpAmberBorder}{RGB}{196,118,18}
\definecolor{lpGray}{RGB}{96,104,114}
\definecolor{lpGrayLight}{RGB}{238,240,243}
\definecolor{lpGreen}{RGB}{34,120,72}
\definecolor{lpGreenLight}{RGB}{222,238,228}
\definecolor{lpGreenBorder}{RGB}{24,96,58}

\definecolor{lpL1}{RGB}{34,120,72}
\definecolor{lpL2}{RGB}{22,90,140}
\definecolor{lpL3}{RGB}{170,42,52}

\usepackage[most]{tcolorbox}
\tcbuselibrary{breakable,skins}

\newtcolorbox{infobox}[1][Hinweis]{%
  enhanced, breakable,
  colback=lpAccentLight,
  colframe=lpAccent,
  boxrule=0.6pt,
  arc=2pt,
  left=10pt,right=10pt,top=8pt,bottom=8pt,
  fonttitle=\sffamily\bfseries\color{white},
  coltitle=white,
  title={#1},
  attach boxed title to top left={xshift=8pt,yshift=-8pt},
  boxed title style={size=small, colback=lpAccent, colframe=lpAccent, boxrule=0.4pt, arc=1pt},
  top=14pt
}

\newtcolorbox{loesungbox}[1][Musterl\"osung]{%
  enhanced, breakable,
  colback=lpGreenLight,
  colframe=lpGreenBorder,
  boxrule=0.6pt,
  arc=2pt,
  left=10pt,right=10pt,top=8pt,bottom=8pt,
  fonttitle=\sffamily\bfseries\color{white},
  coltitle=white,
  title={#1},
  attach boxed title to top left={xshift=8pt,yshift=-8pt},
  boxed title style={size=small, colback=lpGreenBorder, colframe=lpGreenBorder, boxrule=0.4pt, arc=1pt},
  top=14pt
}

\usepackage{enumitem}
\setlist{itemsep=2pt, topsep=4pt, parsep=0pt}
\setlist[itemize,1]{label={\textbullet},leftmargin=1.6em}
\setlist[itemize,2]{label={\textendash},leftmargin=1.4em}
\setlist[enumerate,1]{label=\arabic*., leftmargin=1.8em}

\usepackage{tabularx}
\usepackage{array}
\usepackage{booktabs}
\usepackage{longtable}

\usepackage{fancyhdr}
\usepackage{lastpage}
\pagestyle{fancy}
\fancyhf{}
\renewcommand{\headrulewidth}{0.3pt}
\renewcommand{\footrulewidth}{0pt}
\fancyhead[L]{\footnotesize\sffamily\color{lpGray}L\"osungsheft \textperiodcentered{} Kurs 71 \textperiodcentered{} Tag 2}
\fancyhead[R]{\footnotesize\sffamily\color{lpGray}UML \textperiodcentered{} VSM \textperiodcentered{} Process Mining}
\fancyfoot[L]{\footnotesize\sffamily\color{lpGray}\textcopyright{} Can Siebert}
\fancyfoot[R]{\footnotesize\sffamily\color{lpGray}\thepage{} / \pageref{LastPage}}

\fancypagestyle{plain}{%
  \fancyhf{}%
  \renewcommand{\headrulewidth}{0pt}%
  \fancyfoot[C]{\footnotesize\sffamily\color{lpGray}\thepage{} / \pageref{LastPage}}%
}

\usepackage[explicit]{titlesec}

\titleformat{\section}[block]
  {\sffamily\Large\bfseries\color{lpAccent}}
  {\thesection.}{0.6em}{#1}
  [{\vspace{2pt}\color{lpAccent}\titlerule[0.6pt]}]

\titleformat{\subsection}[block]
  {\sffamily\large\bfseries\color{lpInk}}
  {\thesubsection}{0.6em}{#1}

\titlespacing*{\section}{0pt}{14pt}{8pt}
\titlespacing*{\subsection}{0pt}{10pt}{4pt}

\usepackage[hidelinks,unicode]{hyperref}

\newcommand{\akzent}[1]{{\caveatfont\color{lpAmber}#1}}
\newcommand{\fachbegriff}[1]{\textbf{#1}}
\newcommand{\quelle}[1]{{\footnotesize\sffamily\color{lpGray}(#1)}}

\newcommand{\niveaubadge}[2]{%
  \tcbox[on line, colback=#1, colframe=#1, coltext=white,
         boxrule=0pt, arc=2pt, left=5pt,right=5pt,top=1pt,bottom=1pt,
         nobeforeafter]{\sffamily\bfseries\footnotesize #2}%
}
\newcommand{\nivLi}{\niveaubadge{lpL1}{L1\,\textbullet\,Wissen}}
\newcommand{\nivLii}{\niveaubadge{lpL2}{L2\,\textbullet\,Anwendung}}
\newcommand{\nivLiii}{\niveaubadge{lpL3}{L3\,\textbullet\,Management}}

\newcommand{\zeit}[1]{%
  \tcbox[on line, colback=lpGrayLight, colframe=lpGrayLight, coltext=lpInk,
         boxrule=0pt, arc=2pt, left=5pt,right=5pt,top=1pt,bottom=1pt,
         nobeforeafter]{\sffamily\footnotesize\textbf{$\sim$\,#1}}%
}

\newcommand{\aufgabenkopf}[4]{%
  \par\vspace{6pt}
  \noindent{\sffamily\Large\bfseries\color{lpAccent}Aufgabe #1 \textendash{} #2}\par
  \vspace{2pt}
  \noindent{\color{lpAccent}\rule{\linewidth}{0.6pt}}\par
  \vspace{4pt}
  \noindent #3 \hfill \zeit{#4}\par
  \vspace{6pt}
}

\newcommand{\ytquery}[1]{%
  \par\vspace{4pt}
  \noindent{\footnotesize\sffamily\color{lpGray}\textbf{YouTube-Suchquery:} \texttt{#1}}\par
}

\begin{document}

\begin{titlepage}
  \pagestyle{empty}
  \vspace*{1.5cm}
  \noindent{\sffamily\footnotesize\color{lpGray}LERNPLATT \textperiodcentered{} BY CAN SIEBERT}\\[2pt]
  \noindent{\color{lpAccent}\rule{\linewidth}{1.2pt}}\\[4pt]
  \noindent{\sffamily\footnotesize\color{lpGray}L\"OSUNGSHEFT \textperiodcentered{} MODUL 1 \textperiodcentered{} TAG 2 VON 3 \textperiodcentered{} KURS 71 (Dig 42) \textperiodcentered{} 29.\,Mai 2026}

  \vspace{3cm}
  {\sffamily\fontsize{30}{34}\selectfont\bfseries\color{lpInk}L\"osungsheft\\Tag 2\par}

  \vspace{0.7cm}
  {\sffamily\large\color{lpGray}Musterl\"osungen zu den elf Aufgaben\\des Aufgabenhefts \textendash{} UML, VSM\\und Process Mining im Dreiklang.\par}

  \vspace{1.5cm}
  {\caveatfont\LARGE\color{lpGreen}Musterl\"osungen \textperiodcentered{} nicht die einzige Antwort}\par

  \vfill

  \noindent\begin{minipage}{0.55\linewidth}
    {\sffamily\footnotesize\color{lpGray}Hinweis:}\\
    {\sffamily\small\color{lpInk}Musterl\"osungen sind Diskussionsbasis,\\nicht die einzige korrekte L\"osung.}\\[10pt]
    {\sffamily\footnotesize\color{lpGray}Begleitend:}\\
    {\sffamily\small\color{lpInk}Aufgabenheft \& Lehrskript Tag 2}
  \end{minipage}\hfill
  \begin{minipage}{0.4\linewidth}
    \raggedleft
    {\sffamily\footnotesize\color{lpGray}Dozent:}\\
    {\sffamily\small\color{lpInk}Can Siebert}\\[10pt]
    {\sffamily\footnotesize\color{lpGray}Niveau-Mix:}\\
    {\sffamily\small\color{lpInk}3\,L1 \textperiodcentered{} 5\,L2 \textperiodcentered{} 3\,L3}
  \end{minipage}

  \vspace{0.5cm}
  \noindent{\color{lpAccent}\rule{\linewidth}{0.6pt}}
\end{titlepage}

\section*{Hinweise zu den Musterl\"osungen}
\thispagestyle{plain}

\noindent Die folgenden Musterl\"osungen sind \emph{eine} m\"ogliche, didaktisch nachvollziehbare L\"osung. Insbesondere auf L2 und L3 sind mehrere Begr\"undungswege legitim, solange sie:

\begin{itemize}
  \item den Methodenapparat (UML, VSM, Event-Log-Logik) sauber nutzen,
  \item Annahmen explizit machen, statt sie zu verschleiern,
  \item Zielkonflikte (z.\,B.\ Speed vs.\ Compliance, Daten vs.\ Modell) benennen, statt sie wegzudefinieren,
  \item zu einer klar formulierten Entscheidung f\"uhren \textendash{} nicht blo{\ss} zu einer Beschreibung.
\end{itemize}

\begin{infobox}[Nutzung im Coaching]
Vergleichen Sie Ihre eigene Antwort \emph{vor} der Musterl\"osung. Wo weicht Ihre Argumentation ab? Ist die Abweichung eine andere Annahme \textendash{} oder ein Denkfehler? Genau dort entsteht der Lerneffekt.
\end{infobox}

\clearpage

% =========================================================================
% L1 AUFGABEN
% =========================================================================
\section*{Niveau L1 \textendash{} Wissen \& Reproduktion}
\addcontentsline{toc}{section}{L1 -- Wissen}

% LERNPLATT_HILFSVIDEOS
\section*{Hilfsvideos zu diesem Tag}
\addcontentsline{toc}{section}{Hilfsvideos zu diesem Tag}
\thispagestyle{plain}

\noindent Die folgenden YouTube-Erkl\"arvideos vermitteln die Inhalte, die Sie zur Bearbeitung der Aufgaben ben\"otigen. Klicken Sie auf den Titel, um das Video direkt zu \"offnen; die vollst\"andige URL steht in Klammern.

\begin{description}[style=nextline,leftmargin=18pt,labelindent=0pt,itemsep=8pt]
  \item[1.~UML Use Case Diagram — Wer will was, und wo verläuft der Scope?] \href{https://youtu.be/IWbGedKa8PU}{\textcolor{lpAccent}{https://youtu.be/IWbGedKa8PU}}\\[2pt]
  {\footnotesize\color{lpGray}(URL: \nolinkurl{https://youtu.be/IWbGedKa8PU})}
  \item[2.~UML Activity Diagram mit Swimlanes — Verantwortung sichtbar machen] \href{https://youtu.be/-b5O4gxpKL0}{\textcolor{lpAccent}{https://youtu.be/-b5O4gxpKL0}}\\[2pt]
  {\footnotesize\color{lpGray}(URL: \nolinkurl{https://youtu.be/-b5O4gxpKL0})}
  \item[3.~Wertstromanalyse in der Praxis — VSM Tutorial] \href{https://youtu.be/4GNJxlRuaTs}{\textcolor{lpAccent}{https://youtu.be/4GNJxlRuaTs}}\\[2pt]
  {\footnotesize\color{lpGray}(URL: \nolinkurl{https://youtu.be/4GNJxlRuaTs})}
  \item[4.~Process Mining einfach erklärt — vom Log zum Insight] \href{https://youtu.be/0QOTuw45zhQ}{\textcolor{lpAccent}{https://youtu.be/0QOTuw45zhQ}}\\[2pt]
  {\footnotesize\color{lpGray}(URL: \nolinkurl{https://youtu.be/0QOTuw45zhQ})}
\end{description}
\vspace{0.6em}
\noindent{\footnotesize\color{lpGray}\textit{Tipp:} \"Offnen Sie das Video, lassen Sie es im Hintergrund laufen und arbeiten Sie parallel an der Aufgabe.}

\clearpage

\aufgabenkopf{1}{Drei Methoden \textendash{} drei Perspektiven}{\nivLi}{8\,min}

\ytquery{UML Use Case Activity Diagram VSM Process Mining Vergleich Zweck}

\begin{loesungbox}
\begin{tabularx}{\linewidth}{@{}p{3.0cm}|X|X|X@{}}
\toprule
\textbf{Methode} & \textbf{Hauptzweck} & \textbf{Output} & \textbf{Zentrale Frage} \\
\midrule
UML Use Case  & Scope und Akteure kl\"aren; Verantwortungen sichtbar machen. & Use-Case-Diagramm (Akteure + Use Cases + Beziehungen). & \emph{Wer will was?} \\
\midrule
UML Activity  & Ablauf und Entscheidungen je Rolle (Swimlanes) modellieren. & Activity-Diagramm mit Entscheidungen, Swimlanes. & \emph{Wie l\"auft es ab?} \\
\midrule
VSM           & End-to-End-Wertfluss + Verschwendung sichtbar machen. & Value Stream Map mit BT/WT, Wert/Nicht-Wert, Engpass. & \emph{Wo geht Wert verloren?} \\
\midrule
Process Mining & Ist-Prozess aus Eventlogs ableiten; Varianten \& Abweichungen sehen. & Discovery-Graph, Variantentabelle, Conformance-Report. & \emph{Was passiert wirklich?} \\
\bottomrule
\end{tabularx}

\smallskip
\textbf{Schlusssatz:} F\"ur die Frage \emph{``Wo gehen unsere Durchlaufzeiten verloren?''} ist \fachbegriff{VSM} der prim\"are Hebel \textendash{} es trennt explizit Bearbeitungszeit von Wartezeit. Process Mining \emph{best\"atigt} und \emph{quantifiziert} das Ergebnis aus den Daten; UML w\"are dort sekund\"ar.
\end{loesungbox}

\clearpage

\aufgabenkopf{2}{VSM-Symbolik und Wertschöpfung}{\nivLi}{8\,min}

\ytquery{Value Stream Mapping VSM Symbole Verschwendung Lean Bearbeitungszeit Wartezeit}

\begin{loesungbox}
\textbf{Definitionen + Kategorie:}

\smallskip
\begin{tabularx}{\linewidth}{@{}p{4.0cm}Xc@{}}
\toprule
\textbf{Begriff} & \textbf{Definition} & \textbf{Kategorie} \\
\midrule
Bearbeitungszeit (BT) & Zeit, in der aktiv am Vorgang gearbeitet wird. & meist VA \\
Wartezeit (WT)        & Zeit, in der der Vorgang liegt, ohne Bearbeitung. & NVA \\
WIP                   & Anzahl der Vorg\"ange ``in Arbeit'' an einem Schritt. & Indikator (NVA wenn hoch) \\
Durchlaufzeit (DLZ)   & Gesamtzeit Start bis Ende (BT + WT $\cdot$ alle Schritte). & Ergebnisgr\"o{\ss}e \\
Engpass               & Schritt mit l\"angster Bearbeitungs- oder Wartezeit; begrenzt den Durchsatz. & meist NVA-Treiber \\
Informationsfluss     & Welche Daten/Signale steuern den Prozess (oft Pfeil oben). & NNVA \\
Materialfluss         & Was \emph{flie{\ss}t} physisch oder digital (Ware, Vorgang). & VA, wenn am Output \\
Queue / Warteschlange & Anzahl wartender Vorg\"ange vor einem Schritt. & NVA-Indikator \\
\bottomrule
\end{tabularx}

\smallskip
\textbf{Beispiele:}
\begin{itemize}
  \item \emph{NNVA (notwendig, aber nicht wertsch\"opfend):} Qualit\"atspr\"ufung in einer regulierten Branche \textendash{} der Kunde zahlt nicht direkt daf\"ur, aber ohne sie ist Lieferung unzul\"assig.
  \item \emph{NVA (Verschwendung):} Mehrfache R\"uckfrage wegen fehlender Pflichtfelder \textendash{} reines Suchen, Warten, Erneut-Erfassen ohne neuen Wert.
\end{itemize}

\emph{Merkregel:} Wertsch\"opfend ist nur, wof\"ur der Kunde bezahlen w\"urde \textendash{} alles andere ist NNVA (regulatorisch / qualit\"atssichernd) oder NVA (Verschwendung).
\end{loesungbox}

\clearpage

\aufgabenkopf{3}{Event Log lesen \textendash{} Pflichtfelder und Stolpersteine}{\nivLi}{10\,min}

\ytquery{Process Mining Event Log Case ID Timestamp Activity Discovery Datenqualitaet}

\begin{loesungbox}
\textbf{1. Pflichtfelder:}
\begin{itemize}
  \item \fachbegriff{Case ID} \textendash{} die Vorgangs- bzw.\ Bestellnummer. Beispiel: \emph{Case A / B / C}.
  \item \fachbegriff{Aktivit\"at} \textendash{} der ausgef\"uhrte Schritt. Beispiel: \emph{Angaben\_pr\"ufen}, \emph{Versand}.
  \item \fachbegriff{Timestamp} \textendash{} der Zeitstempel. Beispiel: \emph{09:00, 09:05}.
\end{itemize}
\emph{Optional, aber sehr n\"utzlich:} Ressource (wer hat es getan), Kosten, Standort, Kanal.

\smallskip
\textbf{2. Zwei Prozessvarianten:}
\begin{itemize}
  \item \emph{Variante 1 (Happy Path):} \emph{Case B} \textendash{} Bestellung \textrightarrow{} Pr\"ufen \textrightarrow{} Zahlung \textrightarrow{} Kommissionieren \textrightarrow{} Versand. Linearer, kurzer Pfad.
  \item \emph{Variante 2 (Nachfrage- bzw.\ Nachschub-Pfad):} \emph{Case A} mit Nachfrage-Schleife (Angaben unvollst\"andig) bzw.\ \emph{Case C} mit Warten\_auf\_Nachschub (Lager nicht verf\"ugbar).
\end{itemize}

\smallskip
\textbf{3. Wahrscheinlichster Engpass:}
\emph{Warten auf Nachschub} (Case C) und \emph{Nachfrage-Schleife} (Case A) \textendash{} jeweils Wartezeiten in Stunden bis Tagen. Im Vergleich zur Bearbeitungszeit (Minuten) dominieren beide klar die Durchlaufzeit. Faustregel: Schauen Sie zuerst auf die gr\"o{\ss}ten Zeitspr\"unge zwischen aufeinander folgenden Timestamps.

\smallskip
\textbf{4. Stolpersteine:}
\begin{enumerate}
  \item \emph{Falsche Case ID:} Zwei verschiedene Vorg\"ange teilen sich aus Versehen eine ID \textrightarrow{} Discovery erfindet einen Pfad, den es gar nicht gibt.
  \item \emph{Zeitstempel-Probleme:} Activity B startet \emph{vor} Activity A \textendash{} entweder Datenfehler oder die Schritte laufen parallel und das Log spiegelt das nicht. Erkennen: $\Delta t < 0$ oder unplausible Reihenfolge.
  \item \emph{Inkonsistente Aktivit\"atsnamen:} \emph{``Angaben pr\"ufen''} vs.\ \emph{``Angaben\_pr\"ufen''} vs.\ \emph{``Pr\"ufung''} \textrightarrow{} Discovery zeigt drei Schritte, wo nur einer ist. Erkennen: Vergleich der Aktivit\"atsliste mit der Fachseite.
\end{enumerate}
\end{loesungbox}

\clearpage

% =========================================================================
% L2 AUFGABEN
% =========================================================================
\section*{Niveau L2 \textendash{} Anwendung \& Transfer}
\addcontentsline{toc}{section}{L2 -- Anwendung}

\aufgabenkopf{4}{\emph{RetailNova} \textendash{} VSM aus Alltagstext}{\nivLii}{25\,min}

\ytquery{Value Stream Mapping E2E Online Shop Bestellung Bearbeitungszeit Wartezeit Engpass}

\begin{loesungbox}
\textbf{1. Scope:} \emph{Start:} Kunde l\"ost Bestellung aus. \emph{Ende:} Ware ist beim Kunden bzw.\ Retoure ist erstattet. Retoure wird als \emph{Nebenpfad} mitmodelliert.

\smallskip
\textbf{2. VSM-Schritte (10 Hauptschritte + 1 Nebenpfad):}

\smallskip
\begin{tabularx}{\linewidth}{@{}clXcc@{}}
\toprule
\# & \textbf{Schritt} & \textbf{Beschreibung} & \textbf{Kategorie} & \textbf{BT / WT} \\
\midrule
1  & Bestellung erfassen          & Webshop \"ubernimmt Daten.                  & NNVA & 1 min / 0 \\
2  & Angaben pr\"ufen             & Daten vollst\"andig?                        & NNVA & 5 min / 2\,h \\
3  & R\"uckfrage Kundendaten      & Schleife bei Unvollst\"andigkeit.            & NVA  & 10 min / 4\,h \\
4  & Zahlung pr\"ufen             & PSP-Antwort.                                & NNVA & 1 min / 5 min \\
5  & Lagerverf\"ugbarkeit pr\"ufen & Auf Nachschub warten m\"oglich.             & NNVA & 1 min / bis 2\,T \\
6  & Kommissionieren              & Pickprozess.                                & VA   & 10 min / 1\,h \\
7  & Verpacken                    & Karton, Label.                              & VA   & 5 min / 0 \\
8  & Versand \"ubergeben          & DHL/UPS-\"Ubergabe.                         & NNVA & 5 min / 12\,h \\
9  & Adress\"anderung (Nebenpfad) & Kunde ruft an.                              & NVA  & 5 min / 0 \\
10 & Retoure pr\"ufen             & Eingang Lager.                              & NNVA & 10 min / 2\,T \\
11 & Erstattung anlegen           & PSP-R\"uckbuchung.                          & NNVA & 5 min / 1\,T \\
\bottomrule
\end{tabularx}

\smallskip
\textbf{3. Engpasshypothese:} \emph{Warten auf Nachschub} (bis 2 Tage WT) und \emph{R\"uckfrage Kundendaten} (4\,h WT) \textendash{} hier liegt die Durchlaufzeit, nicht in der Bearbeitung.

\smallskip
\textbf{4. F\"unf Stichpunkte:}
\begin{itemize}
  \item Top-3 Verschwendungen: R\"uckfrage-Schleife, Warten auf Nachschub, Retoure-Lagerstaus.
  \item Gr\"o{\ss}ter Engpass: Lagerverf\"ugbarkeit (Nachschub) \textendash{} dominiert WT.
  \item Offene Frage an den Fachbereich: \emph{Warum entstehen R\"uckfragen \textendash{} fehlt das Pflichtfeld oder die UX-Pflicht?}
  \item Datenvorschlag: Pflichtfeld \emph{Lagerstatus} in Echtzeit ins Webshop-Frontend.
  \item Quick-Win: Pre-Order-Logik f\"ur kritische SKUs (Verbindlichkeit ohne Wartepfad).
\end{itemize}
\end{loesungbox}

\clearpage

\aufgabenkopf{5}{Datenblick \textendash{} Process-Mining-Logik aus Eventlog}{\nivLii}{20\,min}

\ytquery{Process Mining Discovery Variants Bottleneck Eventlog Analyse Beispiel}

\begin{loesungbox}
\textbf{1. Zwei Prozessvarianten:}
\begin{itemize}
  \item \emph{Variante A \textendash{} Nachfrage-Schleife:} Case A. Pfad: Pr\"ufen \textrightarrow{} Nachfragen \textrightarrow{} Vollst\"andig \textrightarrow{} Zahlung \textrightarrow{} Versand. Kennzeichen: Schleife auf \emph{Angaben\_pr\"ufen}, $\Delta t \approx 90$\,min.
  \item \emph{Variante B \textendash{} Happy Path:} Case B. Pfad: Pr\"ufen \textrightarrow{} Zahlung \textrightarrow{} Kommissionieren \textrightarrow{} Versand. Komplettzeit $\approx 3\,$h\,45\,min.
  \item \emph{Variante C (Bonus) \textendash{} Nachschub-Warten:} Case C. Pfad mit \emph{Warten\_auf\_Nachschub} und Pause von 2 Tagen vor Kommissionierung.
\end{itemize}

\smallskip
\textbf{2. Engpass (ausschlie{\ss}lich anhand der Zeiten):}\\
\emph{Warten\_auf\_Nachschub} bei Case C dominiert \textendash{} 2 Tage L\"ucke zwischen 08:26 und Folgeschritt. Selbst die Nachfrage-Schleife (90\,min) verblasst dagegen. Begr\"undung: Die l\"angste Latenz im Log markiert den Engpass; alles andere ist Hintergrundrauschen.

\smallskip
\textbf{3. Drei Hypothesen f\"ur Abweichungen:}
\begin{enumerate}
  \item \emph{Datenqualit\"at:} Fehlende Pflichtfelder im Webshop-Formular (Folge: R\"uckfragen).
  \item \emph{Bestandslage:} Lagerbestand wird zu sp\"at aktualisiert, Bestellung wird angenommen \emph{obwohl} Lager leer.
  \item \emph{Kundeninteraktion:} Telefonische Adress\"anderungen unterbrechen den Standardprozess und werden nicht im Log gef\"uhrt.
\end{enumerate}

\smallskip
\textbf{4. Zwei fehlende Datenfelder:}
\begin{itemize}
  \item \fachbegriff{Ressource} \textendash{} \emph{Wer} hat die Aktivit\"at ausgef\"uhrt? Erlaubt Engpassanalyse pro Person/Team.
  \item \fachbegriff{Lagerstatus / Grundcode} \textendash{} \emph{Warum} Nachschub gewartet wurde (Lieferanten-Verzug, Out-of-Stock, Forecast-Fehler).
\end{itemize}
\end{loesungbox}

\clearpage

\aufgabenkopf{6}{UML Use Case \textendash{} \emph{CloudWorks B2B} Onboarding}{\nivLii}{22\,min}

\ytquery{UML Use Case Diagram Beispiel Customer Onboarding Akteure SaaS}

\begin{loesungbox}
\textbf{1. Use-Case-Set (Textform; Diagramm mit 6 Akteuren, 8 Use Cases):}

\smallskip
\textbf{System:} \emph{CloudWorks Onboarding-Plattform}.\\
\textbf{Akteure:} Kunde \textperiodcentered{} Sales \textperiodcentered{} Onboarding-Team \textperiodcentered{} IT \textperiodcentered{} Compliance \textperiodcentered{} Support.

\smallskip
\textbf{Use Cases + initiierende Akteure:}
\begin{enumerate}
  \item \emph{Account anlegen} \textendash{} Sales (initiiert), Onboarding (f\"uhrt aus).
  \item \emph{Vertrags-/Scope-Check} \textendash{} Sales + Compliance.
  \item \emph{Zug\"ange vergeben} \textendash{} IT (f\"uhrt aus), Onboarding (initiiert).
  \item \emph{Datenimport} \textendash{} Kunde liefert, IT importiert.
  \item \emph{Security / Compliance-Check} \textendash{} Compliance.
  \item \emph{Kunden-Training} \textendash{} Onboarding f\"uhrt mit Kunde durch.
  \item \emph{Go-Live-Freigabe} \textendash{} Onboarding + Compliance.
  \item \emph{Hypercare starten} \textendash{} Support \"ubernimmt von Onboarding.
\end{enumerate}

\smallskip
\textbf{2. Beziehungen:}
\begin{itemize}
  \item Sales \textrightarrow{} \emph{Account anlegen}, \emph{Vertrags-/Scope-Check}.
  \item Onboarding-Team \textrightarrow{} \emph{Datenimport koordinieren}, \emph{Kunden-Training}, \emph{Go-Live-Freigabe}.
  \item IT \textrightarrow{} \emph{Zug\"ange vergeben}, \emph{Datenimport (technisch)}.
  \item Compliance \textrightarrow{} \emph{Security-Check}, \emph{Vertrags-Check}, \emph{Go-Live-Freigabe}.
  \item Support \textrightarrow{} \emph{Hypercare starten}.
\end{itemize}

\smallskip
\textbf{3. Out-of-Scope (6 Wochen):}
\begin{itemize}
  \item \emph{Migration historischer Daten} \textendash{} hoher Aufwand, nicht in 6\,Wochen mit 60\,k umsetzbar.
  \item \emph{Custom-Integrationen je Kunde} \textendash{} z\"ahlt zum erweiterten Hypercare, nicht zum MVP-Onboarding.
\end{itemize}

\smallskip
\textbf{4. Zentrale Verantwortung pro Akteur:}
\begin{itemize}
  \item \emph{Kunde:} liefert vollst\"andige Stammdaten.
  \item \emph{Sales:} \"ubergibt sauberes Briefing (``Definition of Ready'').
  \item \emph{Onboarding:} steuert E2E-Durchlauf bis Go-Live.
  \item \emph{IT:} stellt Zug\"ange und Datenimport in SLA bereit.
  \item \emph{Compliance:} Gate-Entscheidung Security-/Vertrags-Check.
  \item \emph{Support:} \"Ubernahme im Hypercare, nicht vorher.
\end{itemize}
\end{loesungbox}

\clearpage

\aufgabenkopf{7}{UML Activity \textendash{} Kernablauf mit Swimlanes}{\nivLii}{25\,min}

\ytquery{UML Activity Diagram Swimlanes Beispiel Onboarding Entscheidung Eskalation}

\begin{loesungbox}
\textbf{Soll-Logik (textuell):}

\smallskip
\textbf{Swimlanes:} \emph{Sales} \textperiodcentered{} \emph{Onboarding} \textperiodcentered{} \emph{IT} \textperiodcentered{} \emph{Compliance}.

\smallskip
\textbf{Start} (Sales): \emph{Onboarding-Briefing \"ubergeben}.
\begin{itemize}
  \item Activity (Onboarding): \emph{Kundendaten konsolidieren}.
  \item \textbf{XOR}: \emph{Kundendaten vollst\"andig?}
  \begin{itemize}
    \item Nein \textrightarrow{} Activity (Onboarding): \emph{Daten nachfordern} \textrightarrow{} zur\"uck zu \emph{Daten konsolidieren} (max.\ 1 Schleife).
    \item Ja \textrightarrow{} weiter.
  \end{itemize}
  \item Activity (IT): \emph{Account provisionieren}.
  \item \textbf{XOR}: \emph{Provisioning abgeschlossen?}
  \begin{itemize}
    \item Nein nach 48\,h \textrightarrow{} \emph{Eskalation an IT-Lead} (Timer-Trigger).
    \item Ja \textrightarrow{} weiter.
  \end{itemize}
  \item Activity (Compliance): \emph{Security-/Compliance-Check}.
  \item \textbf{XOR}: \emph{Security Check bestanden?}
  \begin{itemize}
    \item Nein \textrightarrow{} Activity (IT/Onboarding): \emph{Korrekturen umsetzen} \textrightarrow{} zur\"uck zum Check.
    \item Ja \textrightarrow{} Activity (Onboarding): \emph{Kunden-Training durchf\"uhren} \textrightarrow{} Activity (Onboarding): \emph{Go-Live-Freigabe}.
  \end{itemize}
\end{itemize}
\textbf{Ende:} \emph{Go-Live abgeschlossen}; \"Ubergang zu Hypercare (Support).

\smallskip
\textbf{Modellzweck:} ``Operatives Steuerungsmodell f\"ur den Onboarding-Lead \textendash{} klare Verantwortung pro Lane und ein Eskalations\-trigger nach 48\,h.''

\smallskip
\textbf{Drei Annahmen:}
\begin{enumerate}
  \item Daten\-nachforderung darf nur \emph{einmal} pro Vorgang laufen; danach Eskalation an Sales.
  \item Provisioning ist technisch automatisierbar, sobald die Pflichtfelder geliefert sind.
  \item Compliance kann innerhalb von 24\,h reagieren, sobald die Pakete vorliegen.
\end{enumerate}

\smallskip
\textbf{Zwei Fragen an den Fachbereich:}
\begin{enumerate}
  \item Welche \emph{Pflichtfelder} sind im Briefing tats\"achlich nicht verhandelbar?
  \item Welcher Compliance-Check ist \emph{vorab} parallelisierbar (vor Provisioning), welcher zwingend sequentiell?
\end{enumerate}
\end{loesungbox}

\clearpage

\aufgabenkopf{8}{Goldene Klammer \textendash{} drei Sichten auf eine Story}{\nivLii}{22\,min}

\ytquery{End to End Process Methodik UML VSM Process Mining Triangulation Onboarding}

\begin{loesungbox}
\textbf{1. UML-Sicht (Struktur):}
Die unscharfe Verantwortung liegt bei \emph{Sales \textrightarrow{} Onboarding}: 40\,\% R\"uckspr\"unge wegen ``fehlender Kundendaten'' bedeuten, dass die \emph{\"Ubergabe} keine sauberen Pflichtfelder erzwingt. Im Activity-Diagramm muss eine Swimlane ``Sales'' eine Aktivit\"at \emph{``Definition of Ready abnehmen''} bekommen, mit einer Entscheidung \emph{``DoR erf\"ullt?''}. Heute ist diese Verantwortung verteilt \textendash{} sie muss \emph{eindeutig} bei Sales liegen.

\smallskip
\textbf{2. VSM-Sicht (Wertfluss):}
Der dominante Engpass ist nicht die 40\,\%-Quote allein, sondern die \emph{absolute Wartezeit}: 25\,\% der F\"alle warten \emph{$>$\,5 Tage} auf IT-Provisioning. R\"uckspr\"unge kosten Stunden, Provisioning-Wartezeit kostet Tage. Damit dominiert \emph{IT-Provisioning} die Durchlaufzeit \textendash{} hier sitzt der gr\"o{\ss}te Hebel.

\smallskip
\textbf{3. Process-Mining-Sicht (Datenrealit\"at):}
\emph{Best\"atigt:} Die VSM-Hypothese ``IT-Wartezeit dominiert'' wird durch die 25-\%-Quote gest\"utzt. \emph{Widerlegt} oder zumindest \emph{relativiert:} Die UML-Hypothese ``Onboarding ist der Engpass'' \textendash{} die Daten zeigen, dass das Onboarding-Team \emph{wartet}, nicht \emph{arbeitet}. 15\,\% Support-Eskalationen vor Go-Live sind ein Sp\"atsignal: das Modell muss eine Support-Lane oder zumindest einen Eskalations-Trigger sichtbar einbauen.

\smallskip
\textbf{Entscheidung:}\\
Wir greifen zuerst beim \fachbegriff{IT-Provisioning} ein \textendash{} weil dort der gr\"o{\ss}te Zeitanteil verloren geht und eine SLA-Eskalationsregel \emph{ohne neues Tooling} umsetzbar ist (48-h-Trigger, Standard\-pakete f\"ur Zug\"ange). Die Daten-R\"uckspr\"unge werden in Phase 2 \"uber eine Definition of Ready adressiert \textendash{} der Hebel ist kleiner, das Ergebnis aber nachhaltiger.
\end{loesungbox}

\clearpage

% =========================================================================
% L3 AUFGABEN
% =========================================================================
\section*{Niveau L3 \textendash{} Management \& Entscheidungsarchitektur}
\addcontentsline{toc}{section}{L3 -- Management}

\aufgabenkopf{9}{Fallstudie \emph{CloudWorks B2B}}{\nivLiii}{45\,min}

\ytquery{Customer Onboarding SaaS Operating Model VSM UML Process Mining Fallstudie}

\begin{loesungbox}
\textbf{1. Use-Case-Set:} Account anlegen \textperiodcentered{} Vertrags-/Scope-Check \textperiodcentered{} Zug\"ange vergeben \textperiodcentered{} Datenimport \textperiodcentered{} Security/Compliance-Check \textperiodcentered{} Kunden-Training \textperiodcentered{} Go-Live-Freigabe \textperiodcentered{} Hypercare starten.

\smallskip
\textbf{2. Activity-Kernlogik (textform):}\\
\emph{Start} \textrightarrow{} Sales \"ubergibt Onboarding-Briefing \textrightarrow{} \textbf{XOR} ``Kundendaten vollst\"andig?'' (Nein \textrightarrow{} Daten nachfordern \textrightarrow{} zur\"uck; max.\ 1 Schleife) / (Ja \textrightarrow{} IT provisioniert) \textrightarrow{} \textbf{XOR} ``Provisioning abgeschlossen?'' (Nein nach 48\,h \textrightarrow{} Eskalation) \textrightarrow{} Compliance-Check \textrightarrow{} \textbf{XOR} ``Bestanden?'' (Nein \textrightarrow{} Korrekturen) \textrightarrow{} Training \textrightarrow{} Go-Live-Freigabe \textrightarrow{} \emph{Ende}.

\smallskip
\textbf{3. VSM (E2E, 10 Schritte):}\\
Briefing \textrightarrow{} Daten konsolidieren (NNVA, 1\,h BT / 4\,h WT) \textrightarrow{} R\"uckfrage (NVA, 30\,min / 1\,T) \textrightarrow{} Provisioning (NNVA, 30\,min / 5\,T \textendash{} \fachbegriff{Engpass}) \textrightarrow{} Konfiguration (VA, 4\,h / 0) \textrightarrow{} Security-Check (NNVA, 2\,h / 1\,T) \textrightarrow{} Korrekturen (NVA, 2\,h / 1\,T) \textrightarrow{} Training (VA, 4\,h / 0) \textrightarrow{} Go-Live-Freigabe (NNVA, 30\,min / 1\,T) \textrightarrow{} Hypercare (VA, n.\,d.).

\smallskip
\textbf{4. Process-Mining-Hypothesen:}
\begin{enumerate}
  \item Die 25\,\% Provisioning-Wartezeit korrelieren mit \emph{Kundengr\"o{\ss}e} (gr\"o{\ss}ere Kunden brauchen mehr Standard\-pakete) \textendash{} Hypothese pr\"ufen, sobald Kunden-Tier im Log.
  \item Die 40\,\% R\"uckspr\"unge konzentrieren sich auf \emph{wenige Sales-Mitarbeitende} \textendash{} Coaching-Hebel, kein Tool-Problem.
\end{enumerate}

\smallskip
\textbf{5. Zwei Ma{\ss}nahmen in 6 Wochen:}
\begin{enumerate}
  \item \fachbegriff{Definition of Ready (DoR)} f\"ur Onboarding: Pflichtfelder + \"Ubergabeformat. \emph{Impact:} reduziert R\"uckspr\"unge sofort. \emph{Risiko:} Sales-Widerstand. \emph{Aufwand:} 1\,Person, 2\,Wochen.
  \item \fachbegriff{SLA-Eskalationsregel} f\"ur Provisioning (48\,h-Trigger) + 3 Standard\-pakete f\"ur Zug\"ange. \emph{Impact:} reduziert Wartezeit drastisch. \emph{Risiko:} IT-Kapazit\"atsspitzen. \emph{Aufwand:} 1\,Person IT + 0,5\,Person PMO, 3\,Wochen.
\end{enumerate}
\emph{Beide ohne neues Tool umsetzbar} \textendash{} passt in 60-k-Budget.

\smallskip
\textbf{6. Governance-Mini:}
\begin{itemize}
  \item Onboarding-Checkliste: Owner = Head of Onboarding.
  \item IT-SLA: Owner = IT Operations Lead.
  \item Compliance-Gate: Owner = Compliance Officer.
  \item KPIs: \emph{Durchlaufzeit} (Median), \emph{R\"ucksprung-Rate} (\% F\"alle mit Datenr\"ucksprung), \emph{Eskalationsquote} (\% F\"alle mit 48-h-Trigger).
\end{itemize}
\end{loesungbox}

\clearpage

\aufgabenkopf{10}{Methodenauswahl \textendash{} Wann was?}{\nivLiii}{30\,min}

\ytquery{Methodenauswahl Process Mining VSM UML Reifegrad Operating Model Entscheidung}

\begin{loesungbox}
\textbf{1. Methodenlandkarte:}

\smallskip
\begin{tabularx}{\linewidth}{@{}p{2.4cm}XXXX@{}}
\toprule
\textbf{Methode} & \textbf{Zweck} & \textbf{Zielgruppe} & \textbf{Output} & \textbf{Qualit\"ats\-schwelle} \\
\midrule
UML Use Case   & Scope und Akteure kl\"aren. & Sales, Produkt, Fachbereich. & Use-Case-Diagramm. & 6\,--\,8 Use Cases, alle Akteure benannt. \\
\midrule
UML Activity   & Ablauf, Entscheidungen, Rollen. & Onboarding-Lead, Operations. & Activity-Diagramm mit Swimlanes. & 2\,--\,3 Swimlanes, $\geq$\,2 Entscheidungen. \\
\midrule
VSM            & Wertfluss + Verschwendung. & Operations, COO. & E2E-VSM mit BT/WT, Engpass. & 8\,--\,12 Schritte, Engpass markiert. \\
\midrule
Process Mining & Ist-Prozess + Varianten aus Daten. & Data/Process Owner, Audit. & Discovery-Graph, Variantentabelle. & Mind.\ 3 Pflichtfelder im Log, plausible Varianten. \\
\bottomrule
\end{tabularx}

\smallskip
\textbf{2. Warum kombinieren?}
Jede Methode hat einen blinden Fleck: UML zeigt Verantwortung, aber \emph{nicht} Zeit; VSM zeigt Wertfluss, aber \emph{nicht} Varianten\-vielfalt; Process Mining zeigt Daten, aber \emph{nicht} Intentionen. Die ``goldene Klammer'' bedeutet: eine Story, drei Perspektiven. Eine einzige Methode l\"asst Sie systematisch dort blind, wo Sie am wenigsten erwarten \textendash{} besonders bei der \"Ubergabe zwischen Rollen, an Systemgrenzen und in Ausnahmef\"allen.

\smallskip
\textbf{3. Annahmen:}
\begin{itemize}
  \item \emph{Akzeptabel:} (a) Sch\"atzungen f\"ur BT/WT in VSM (Gr\"o{\ss}enordnung gen\"ugt). (b) Aktivit\"atsnamen werden normiert. (c) Erster Discovery-Lauf ohne Ressource-Feld.
  \item \emph{Risikorelevant (dokumentationspflichtig):} (a) Wir lassen 10\,\% Cases mit unsauberen Case-IDs aus der Analyse weg. (b) Wir setzen ``Provisioning bestanden'' = Timestamp \emph{Account aktiv}, obwohl das nur ein Proxy ist.
\end{itemize}

\smallskip
\textbf{4. Faustregel:}
Process Mining lohnt sich ab dem Reifegrad, in dem \emph{(a)} ein zentrales System mindestens 80\,\% der Aktivit\"aten loggt, \emph{(b)} Case-IDs konsistent vergeben werden und \emph{(c)} die Organisation Abweichungen als \emph{Signal} und nicht als ``Schuldzuweisung'' interpretiert. Vorher: VSM + UML, plus eine manuelle Stichprobe.
\end{loesungbox}

\clearpage

\aufgabenkopf{11}{Transferprojekt \textendash{} E2E-Operating Model}{\nivLiii}{45\,min}

\ytquery{E2E Operating Model Process Mining Governance Decision Architecture Roadmap}

\begin{loesungbox}
\emph{Musterl\"osung als Entscheidungsarchitektur \textendash{} andere Konfigurationen sind m\"oglich.}

\smallskip
\textbf{1. Methodenlandkarte:} siehe Aufgabe 10 (UML / BPMN / EPC / VSM / Mining). Erg\"anzung: \emph{BPMN} f\"ur ausf\"uhrbare Prozesse, \emph{EPC} f\"ur Management-\"Ubersicht in ARIS-Welt.

\smallskip
\textbf{2. Daten-Minimum (Eventlog):}
\begin{itemize}
  \item Pflichtfelder: \emph{Case ID}, \emph{Aktivit\"at}, \emph{Timestamp}, \emph{Ressource}, \emph{Standort}.
  \item Verantwortliche: Data Steward je Quellsystem (CRM, Ticketing, ERP).
  \item Datenquellenstrategie: \emph{primary first} \textendash{} CRM als Master f\"ur Case-IDs; weitere Quellen \"uber konsistenten Customer-Key joinen.
\end{itemize}

\smallskip
\textbf{3. Governance:}
\begin{itemize}
  \item \emph{Process Owner} (Fachbereich) \textendash{} Accountable f\"ur Prozess + Modell.
  \item \emph{Modeler} (BPM-Team) \textendash{} Responsible f\"ur UML/BPMN/VSM-Artefakte.
  \item \emph{Data Steward} \textendash{} Responsible f\"ur Log-Qualit\"at je System.
  \item \emph{Reviewer} (Quality) \textendash{} Vier-Augen-Pr\"ufung.
  \item Review-Zyklus: monatliches Process Board; Eskalation an CTO bei Zielkonflikten Speed vs.\ Compliance.
\end{itemize}

\smallskip
\textbf{4. Entscheidungsarchitektur \textendash{} Top-5 Entscheidungen:}
\begin{enumerate}
  \item \emph{Priorisierung:} Welche Prozesse zuerst harmonisieren (Volumen $\times$ Schmerz)?
  \item \emph{Gate-Freigabe:} Wer entscheidet \"uber Prozess-/Workflow-\"Anderungen?
  \item \emph{Automatisierungs\-kandidaten:} Welche Prozesse sind reif f\"ur BPMS/RPA?
  \item \emph{Compliance-Risiken:} Wo m\"ussen Steuerungs\-punkte sitzen?
  \item \emph{Kapazit\"aten:} Welche Engp\"asse (Person, System) limitieren das Wachstum?
\end{enumerate}

\smallskip
\textbf{5. Roadmap (10 Wochen):}
\begin{itemize}
  \item \emph{Wo 1\,--\,2:} Methodenlandkarte + Pilotprozess (\emph{Customer Onboarding}) festlegen.
  \item \emph{Wo 3\,--\,5:} Use-Case + Activity + VSM des Piloten; Eventlog-Spezifikation.
  \item \emph{Wo 6\,--\,8:} Erster Discovery-Lauf; KPI-Set verproben (DLZ, R\"uckspr\"unge, Varianten).
  \item \emph{Wo 9\,--\,10:} Governance Board konstituieren, Lessons Learned, Rollout-Plan f\"ur den n\"achsten Prozess.
\end{itemize}

\smallskip
\textbf{6. Zwei Entscheidungen unter Unsicherheit:}
\begin{enumerate}
  \item \emph{Pilotprozess = Onboarding} (statt z.\,B.\ Procure-to-Pay). Annahme: hohe Sichtbarkeit + Kunden-Impact. Verworfen: P2P (gro{\ss}e IT-Abh\"angigkeit). \emph{Fr\"uhwarnsignal:} Wenn nach 6 Wochen weniger als 5 Cases sauber durchanalysiert sind \textrightarrow{} Pilot wechseln oder Scope verkleinern.
  \item \emph{Minimal-Standard = nur 3 Pflichtfelder im Log} (statt 8). Annahme: Akzeptanz vor Vollst\"andigkeit. Verworfen: ``Full-Schema'' (bremst Quellsysteme). \emph{Fr\"uhwarnsignal:} Wenn Conformance-Aussagen wegen fehlender Felder unbegr\"undbar werden \textrightarrow{} 2 Felder nachziehen.
\end{enumerate}

\smallskip
\textbf{7. Top-5 Risiken:}
\begin{itemize}
  \item Akzeptanz im Fachbereich \textrightarrow{} Key-User fr\"uh einbeziehen.
  \item Datenqualit\"at \textrightarrow{} Data-Quality-Checks im ETL erzwingen.
  \item \"Ubermodellierung \textrightarrow{} ``Modell folgt Zweck''-Pflicht.
  \item Tool-Lock-in \textrightarrow{} XML-Export als Auswahl\-kriterium.
  \item Reorganisation \textrightarrow{} Rollen an Funktion, nicht an Personen koppeln.
\end{itemize}

\smallskip
\textbf{8. Anschluss an Strategie:}
\begin{enumerate}
  \item \emph{Compliance:} sichtbare Prozesse reduzieren Audit-Findings.
  \item \emph{Time-to-Market:} klare Eigent\"umerschaft verk\"urzt Eskalationen.
  \item \emph{Skalierbarkeit:} ein gemeinsames Daten- und Modellverst\"andnis macht Wachstum (neue Standorte, Zuk\"aufe) \"uberhaupt erst beherrschbar.
\end{enumerate}
\end{loesungbox}

\clearpage

\section*{Abschluss}
\thispagestyle{plain}

\noindent Die Musterl\"osungen sind eine Diskussionsbasis. Wenn Ihre eigene L\"osung anders ist, fragen Sie: \emph{Habe ich andere Annahmen \textendash{} oder einen anderen Modellzweck angelegt?} Beides ist legitim, solange die Logik nachvollziehbar bleibt.

\medskip
\begin{infobox}[N\"achster Schritt]
Bringen Sie Ihre eigenen L\"osungen in die Coaching-Sitzung mit. Im Fokus stehen drei Fragen: Welche Annahmen haben Sie gemacht? Welche Zielkonflikte (Speed vs.\ Compliance, Modell vs.\ Daten) haben Sie sichtbar gemacht? Welche Entscheidung haben Sie getroffen \textendash{} und w\"urden Sie auch verantworten, wenn sich die Datenlage \"andert?
\end{infobox}

\vspace{1cm}
\noindent{\color{lpAccent}\rule{\linewidth}{0.6pt}}\\[4pt]
{\footnotesize\sffamily\color{lpGray}Lernplatt \textperiodcentered{} by Can Siebert \textperiodcentered{} Kurs 71 (Dig 42) \textperiodcentered{} Tag 2 \textperiodcentered{} L\"osungsheft \textperiodcentered{} \textcopyright{} 2026}

\end{document}
